\chapter{Probability Logic and Conditional-Event Coherence}
\label{sec:p15-rationality-4-4}

Pfeifer's 2021 chapter presents probability logic as a rationality
framework for reasoning under uncertainty, and compares its formal
standards with experimental response patterns. Its subject is uncertain
arguments and conditional events, not probabilities of temporal paths
in a transition system \cite[pp.~277--284]{Pfeifer2021ProbabilityLogic}.
The treatment below preserves the chapter's three claims: the framework
generalizes bivalent propositional reasoning, admits a specified connexive
interpretation, and blocks unrestricted strengthening of a default's
conditioning context. Mathematical validity and empirical support for
using this framework are distinct claims.

\section{The objects and the consequence question}

\begin{hbdefinition}{def:p15-rat44-events}{Propositions, conditional events, and their possible cases}
Ordinary propositions have classical truth conditions. Following the
chapter, write \(A\mathbin{\&}B\) for conjunction, \(A\lor B\) for
disjunction, \(\sim A\) for negation, and \(\top,\bot\) for a
tautology and contradiction. A finite nonempty constituent space \(\Omega\)
contains the logically possible truth assignments to the propositions
under consideration; an event is a subset of this space.

For \(A\ne\varnothing\), the conditional event \(B\mid A\) is
true on \(A\mathbin{\&}B\), false on
\(A\mathbin{\&}\sim B\), and void on \(\sim A\). A bet on it is
won, lost, or called off, respectively. Its probability is a primitive
assessment, not a third truth value or the probability of the material
conditional \(\sim A\lor B\). The chapter's simple conditional events
have ordinary, nonconditional operands: neither operand contains
another occurrence of \(\mid\).

The event \(A\) may be possible while having probability zero.
This differs from \(A=\bot\), which is excluded as a conditioning
event in this coherence framework. The chapter notes that Popper
functions use a different convention allowing contradictory conditions;
that convention is not inserted here.
\end{hbdefinition}

\begin{hbdefinition}{def:p15-rat44-coherence}{Two-sided conditional betting coherence}
For a finite family \(E_i\mid H_i\), with every \(H_i\ne\varnothing\),
and prices \(p_i\in[0,1]\), define the gain from stakes \(s_i\in\mathbb R\)
by
\[
 G_s(\omega)=\sum_i s_i\mathbf1_{H_i}(\omega)
                            (\mathbf1_{E_i}(\omega)-p_i).
\]
It is zero for an individual called-off bet. On the active event
\(H_s=\bigcup_{i:s_i\ne0}H_i\), coherence requires
\[
             \min_{H_s}G_s\leq0\leq\max_{H_s}G_s
\]
for every nonzero portfolio. An arbitrary assessment family is coherent
when every finite subfamily has this property. Positive and negative
stakes represent buying and selling at the same precise price. This is
a formal no-sure-gain/no-sure-loss condition, not a claim that every
person actually accepts every monetary gamble.

Unconditional events are \(E\mid\top\), so
\(p(E)\defeq p(E\mid\top)\). For specified premise values or
interval constraints, the probability-logical conclusion is the set of
values extendible to a coherent assessment satisfying those constraints.
Its tight lower and upper bounds are the infimum and supremum of this
set, when it is nonempty. An inconsistent premise assessment has no
coherent extension and is reported as such, not as a vacuous interval.
``Noninformative'' means this set is the entire \([0,1]\), not
that the conclusion is false or that the premises lack practical relevance.
\end{hbdefinition}

The constituent constructions and proofs below are finite semantic
reconstructions of the chapter's stated coherence conditions. They do
not claim to transcribe the full external algorithms or books cited
there. Unless a different logical relation is specified, the named
atomic propositions are logically independent: all their truth patterns
are possible. Known logical dependencies can narrow the quoted bounds.
No probabilistic independence assumption is made.

\begin{hbtheorem}{thm:p15-rat44-coherence-method}{proved=annex:p15.rat44.coherence}{Convex geometry and the finite conditional method}
For unconditional events \(E_1,\ldots,E_n\), let
\(Q=\{(\mathbf1_{E_i}(\omega))_{i=1}^n:\omega\in\Omega\}\).
The assessment vector \(p\) is coherent iff \(p\in\operatorname{conv}Q\),
equivalently iff nonnegative constituent weights sum to one and give
the specified event probabilities.

An empty assessment is accepted. For a nonempty finite conditional
assessment family, the procedure is:
\begin{enumerate}
\item On the union of the remaining conditioning events, solve
      \(\lambda_\omega\geq0\), \(\sum\lambda_\omega=1\), and
      \(\sum_\omega\lambda_\omega\mathbf1_{H_i}
                  (\mathbf1_{E_i}-p_i)=0\) for all remaining \(i\).
\item If infeasible, reject the assessment. Otherwise identify the
      conditioning events whose maximum attainable weight over this
      solution set is zero, and repeat on just those assessments.
\item Accept when none remain.
\end{enumerate}
This terminates in at most the number of assessed conditionals and is
equivalent to finite coherence. The repeated zero-weight stages are
essential; one unconditional distribution need not decide a primitive
conditional on a null but possible event.
\end{hbtheorem}

\begin{hbtheorem}{thm:p15-rat44-basic-laws}{proved=annex:p15.rat44.laws}{The basic identities and the positive-probability bridge}
For every admissible \(A\), coherence implies
\[
 p(A\mid A)=1,\quad p(\sim A\mid A)=0,\quad
 p(B\mid A)+p(\sim B\mid A)=1.
\]
On unconditional events it implies the ordinary finite probability
identities. When the joint assessment includes \(p(A)>0\), it implies
\[
              p(B\mid A)=\frac{p(A\mathbin{\&}B)}{p(A)}.
\]
At \(p(A)=0\), the equation
\(p(A\mathbin{\&}B)=p(A)p(B\mid A)\) does not determine the
conditional price. Its admissible values still depend on coherence and
logical relationships on \(A\), rather than on an assigned value for \(0/0\).
\end{hbtheorem}

An unconditional probability function takes values in \([0,1]\), but
is not an arbitrary assignment of graded truth values. The constituent
constraints preserve the underlying Boolean events. For example knowing
\(p(B),p(C)\) generally leaves an interval for
\(p(B\mathbin{\&}C)\). This is the chapter's contrast with
truth-functional connectives: exact marginal beliefs do not determine
the joint belief. Zero-one ordinary assessments recover classical
valuations; void conditional events remain a further, distinct extension.

\section{Zero-probability antecedents and the conditional paradox}

The chapter contrasts ``\(B\); therefore, if \(A\), then \(B\)''
under two representations. As a material conditional the conclusion is
classically entailed. Substituting a green-cheese claim about the moon
for \(A\) and a weather claim about Vienna for \(B\) makes the
pragmatic oddity evident without invalidating the classical argument.

\begin{hbtheorem}{thm:p15-rat44-paradox}{proved=annex:p15.rat44.paradox}{The two tight bounds, including both endpoints}
For logically independent \(A,B\) and every \(x\in[0,1]\),
\[
 \begin{array}{ll}
 p(B)=x:& p(B\mid A)\text{ can be any value in }[0,1];\\
 p(B)=x:& p(\sim A\lor B)\text{ can be any value in }[x,1].
 \end{array}
\]
The first claim includes \(x=0\) and \(x=1\), using primitive
conditional assessments on a possible null antecedent. If instead one
requires \(p(A)>0\) and uses only the quotient definition, then
\(x=0\) forces conditional value zero and \(x=1\) forces value one.
For \(0<x<1\), with no further restriction on positive \(p(A)\),
every conditional value remains attainable.
\end{hbtheorem}

Assigning value one to every conditional on a null antecedent would
assign one to both \(B\mid A\) and \(\sim B\mid A\), violating
the complement identity. The chapter therefore rejects that default
within its coherent-price semantics. This is not an assertion that
every alternative formalism uses its same domain and complement laws.
Nor does \(p(B)=1\) make \(B\) a tautology: if \(B\) is logically
true, coherence does force \(p(B\mid A)=1\) on every admissible \(A\).

\section{Connexivity with a declared negation convention}

\begin{hbdefinition}{def:p15-rat44-connexive-reading}{Narrow-scope conditional negation and high acceptance}
For this part of the chapter, negating ``if \(A\), then \(B\)''
means replacing its consequent by \(\sim B\): the conditional event
is \((\sim B)\mid A\). A high-probability acceptance threshold is
\(h>1/2\), with acceptance requiring \(p(B\mid A)>h\).
The chapter also refers to wide-scope interpretations elsewhere; it
does not identify them with this narrow-scope convention.
The threshold implications below are assertions about assessments,
not nested conditional events of the prohibited form
\(C\mid(B\mid A)\).
\end{hbdefinition}

\begin{hbtheorem}{thm:p15-rat44-connexive}{proved=annex:p15.rat44.connexive}{Aristotle, Boethius, and Abelard in the selected interpretation}
For a contingent \(A\), the Aristotle principles
\(\sim(\text{if }\sim A\text{ then }A)\) and
\(\sim(\text{if }A\text{ then }\sim A)\) are justified by
\[
 p(A\mid\sim A)=p(\sim A\mid A)=0,\qquad
 p(\sim A\mid\sim A)=p(A\mid A)=1.
\]
Under the selected narrow-scope convention, each Boethius principle
\[
 (A\to B)\to\sim(A\to\sim B),\qquad
 (A\to\sim B)\to\sim(A\to B)
\]
preserves its high-probability assertion because double negation gives
the same conditional assessment. Abelard's prohibition of accepting
both \(A\to B\) and \(A\to\sim B\) holds since their conditional
probabilities sum to one and cannot both exceed \(1/2\).

Replacing every arrow by material implication instead makes the two
Aristotle formulas equivalent to \(\sim A,A\), respectively, and
both Boethius formulas and the negation of their opposing conjunction
equivalent to \(A\). They are contingent, not classical tautologies.
The claimed validation is therefore interpretation-specific, not a
proof that these formulas are valid for every meaning of a conditional
or every scope of negation.
\end{hbtheorem}

The chapter reports experimental support for the Aristotle principles
and marks the other connexive principles as needing further empirical
research. It does not transfer the Aristotle result to Boethius or
Abelard merely because all three have a formal justification
\cite[p.~280]{Pfeifer2021ProbabilityLogic}.
Its comparison with connexive logics is not a proposal to add these
formulas, with every arrow still material, to classical logic. The
chapter describes connexive calculi as alternatives not simply contained
in, or proper extensions of, the classical calculus. The assessment
interpretation above is correspondingly not being advertised as a
complete axiomatization of a particular connexive calculus.

\section{Default contexts and tight propagation}

\begin{hbtheorem}{thm:p15-rat44-noninformation}{proved=annex:p15.rat44.noninformation}{Strengthening, contraposition, and ordinary transitivity}
With the indicated propositions logically independent and all
conditioning events possible, the following conclusion values range
over all \([0,1]\), for every stated premise value:
\begin{enumerate}
\item From \(p(F\mid B)=x\): \(p(F\mid(B\mathbin{\&}P))\).
\item From \(p(B\mid A)=x\): \(p(\sim A\mid\sim B)\).
\item From \(p(B\mid A)=x\) and \(p(C\mid B)=y\): \(p(C\mid A)\).
\end{enumerate}
Here \(x,y\in[0,1]\); the endpoint cases can require null-event
conditioning strata. These are the chapter's probabilistically
noninformative counterparts of classically valid material-conditional
arguments, not proofs that every uncertain transitivity inference is
useless once additional dependence information is supplied.
\end{hbtheorem}

The bird/penguin example illustrates why ``normally a bird flies''
does not license the same assessment for a bird known to be a penguin.
``Normally'' scopes over the whole conditional, not only the consequent.
System P is the chapter's reference family for basic default-inference
constraints. It cites Gilio's coherent probability propagation for its
rules and notes that unrestricted strengthening, contraposition, or
transitivity would destroy the intended nonmonotonic discipline
\cite[pp.~280--281, note~5]{Pfeifer2021ProbabilityLogic}.

\begin{hbwarning}{def:p15-rat44-monotonicity-distinction}{Changing a conditioning context is not adding a numerical constraint}
For a fixed target and fixed assessment language, adding numerical
premise constraints can only shrink the set of coherent extensions;
its lower bound cannot decrease and its upper bound cannot increase.
That constraint consequence is monotone. The chapter's nonmonotonicity
concerns revising default conclusions or changing \(B\) to
\(B\mathbin{\&}P\) as the conditioning context. Treating these
two operations as the same would contradict the definition of tight
bounds.
\end{hbwarning}

\begin{hbtheorem}{thm:p15-rat44-propagation}{proved=annex:p15.rat44.propagation}{The selected AND, CUT, and MP bounds}
For \(0\leq x_L\leq x_U\leq1\) and
\(0\leq y_L\leq y_U\leq1\), the constraints
\(p(B\mid A)\in[x_L,x_U]\) and
\(p(C\mid A)\in[y_L,y_U]\) give exactly
\[
 p((B\mathbin{\&}C)\mid A)\in
 [\max\{0,x_L+y_L-1\},\min\{x_U,y_U\}].
\]
For cumulative transitivity, \(p(B\mid A)=x\) and
\(p(C\mid(A\mathbin{\&}B))=y\) give exactly
\[
                         p(C\mid A)\in[xy,xy+1-x].
\]
For probabilistic modus ponens, \(p(B)=x\) and \(p(C\mid B)=y\)
give exactly the same interval for \(p(C)\). All bounds are sharp,
including \(x=0,1\), and do not multiply marginals under an unstated
independence assumption. MP is CUT with \(A=\top\).
\end{hbtheorem}

For example \(p(B\mid A)=0.8,p(C\mid A)=0.7\) gives the AND
interval \([0.5,0.7]\), not the independence point \(0.56\).
MP with \(x=0.8,y=0.7\) instead gives \([0.56,0.76]\).
These calculations use stipulated assessments, not measured responses.
When both AND premise values are one, its conclusion is one; exact
premise values otherwise need not give an exact conclusion value.

\section{What the experimental comparisons establish}

The chapter contrasts the former classical-logic benchmark with the
probabilistic new paradigm. Its examples are Braine and O'Brien's
mental logic and Rips's mental rule theory, using syntactic proof
standards, and Johnson-Laird's mental-model account, using a semantic
model standard. Changing the rationality framework changes what counts
as an error and motivates different tasks; it does not by itself show
that a participant used a particular formal procedure. The chapter's
empirical interest is general response patterns in uncertain reasoning,
not selecting a probability model merely for its fit to data
\cite[pp.~277,279]{Pfeifer2021ProbabilityLogic}.

The ordinary truth-table task asks whether the four cases confirm,
disconfirm, or are irrelevant to an indicative conditional. The chapter
reports common confirmation for \(A\mathbin{\&}B\), disconfirmation
for \(A\mathbin{\&}\sim B\), and irrelevance of \(\sim A\).
Calling this a defective truth table already presupposes the material
conditional as the benchmark; the conditional-event semantics assigns
the reported cases true, false, and void instead.

In the probabilistic truth-table task (PTTT), the input is
\[
 p(A\mathbin{\&}B)=x_1,\quad
 p(A\mathbin{\&}\sim B)=x_2,\quad
 p(\sim A\mathbin{\&}B)=x_3,\quad
 p(\sim A\mathbin{\&}\sim B)=x_4,
\]
where the nonnegative values sum to one. With \(x_1+x_2>0\), the
competing predictions are \(x_1/(x_1+x_2)\) for a conditional event,
\(x_1\) for conjunction, and \(x_1+x_3+x_4\) for a material
conditional. For stipulated \((0.2,0.1,0.3,0.4)\), these give
\(2/3,0.2,0.9\), respectively. They are hypotheses about interpreting
the task, not three estimators of the same defined event. When
\(x_1+x_2=0\), the ratio supplies no answer; a primitive conditional
assessment then requires its own coherent specification rather than
evaluation of an undefined quotient.

The following are the chapter's reports of the cited studies, not fresh
population estimates \cite[pp.~278--281]{Pfeifer2021ProbabilityLogic}:
\begin{itemize}
\item Indicative PTTT tasks, often using abstract shape/color material,
      predominantly elicited conditional-probability responses, with a
      substantial conjunction minority and negligible material-conditional
      responses. Repeated tasks produced a shift toward the conditional
      interpretation; the chapter reports about 80\% of responses in
      the last tasks, citing Fugard et al. (2011). This denominator is
      task responses, not all people or all everyday conditionals.
\item It cites causal, counterfactual, and abductive task studies as
      extending the predictive success of conditional probability. These
      involve different formulations and materials; no common effect
      size or universal equivalence of those tasks is established here.
\item On AND interval tasks it reports predominantly coherent interval
      responses, and mostly lower-bound rather than upper-bound
      violations among the incoherent responses. An upper-bound
      violation would be the formal conjunction fallacy; these errors
      therefore differ from that particular error pattern.
\item It reports support for selected System P intervals and MP, and
      recognition of strengthening and contraposition as noninformative.
      Apparent endorsement of ordinary transitivity is interpreted as
      a possible pragmatic reading as CUT, not directly as evidence
      that the ordinary transitivity bounds become informative.
\end{itemize}

The chapter does not specify sample recruitment, sample sizes, all
response coding decisions, confidence intervals, or task-by-task effect
sizes for this collection of studies. Its broad descriptive conclusion
must be read with that reporting limit; those details are not invented
or pooled here. The Linda-task comparison also raises an interpretation
objection: merely using the words ``and'' and ``probability'' does not
ensure that participants interpreted the ranking task as the formal
conjunction-probability problem. Neither this objection nor the interval
results disprove every finding of a conjunction error. Descriptive fit
alone does not establish the normative authority of coherence.

\section{Nested-conditionals boundary and handbook comparison}

\begin{hbtheorem}{thm:p15-rat44-nested-source}{cited=Pfeifer2021ProbabilityLogic}{The chapter's reported nested-MP extension}
The chapter reports an extension using conditional random quantities
and conditional previsions: from ``the cup breaks if dropped'' and
``if the cup breaks if dropped, then the cup is fragile,'' the lower
conclusion bound is the product of the two premise assessments and
the upper bound is that product plus one minus the first assessment.
This is the reported \([xy,xy+1-x]\) nested-MP result. It requires
the richer conditional-random-quantity framework cited there; it is
not obtained by forming \(p(F\mid(B\mid D))\) in the simple
conditional-event language. The chapter states that import--export
fails in that richer approach and connects the distinction to avoiding
Lewis-type triviality. It gives no complete construction or derivation
of that extension, and explicitly calls for further psychological
evaluation \cite[p.~282, note~1]{Pfeifer2021ProbabilityLogic}.
\end{hbtheorem}

The comparison with existing handbook treatments is exact:
\begin{itemize}
\item Finite probability supplies masses, expectations, and quotient
      conditioning at positive probability
      (\cref{sec:p15-finite-probability,def:p15-finite-conditioning}).
      The primitive conditional assessment used here remains meaningful
      on a possible null event; it does not repair the quotient by
      assigning a value to \(0/0\). Finite-portfolio coherence supplies
      finite additivity; countable additivity in the measurable-probability
      treatment is an additional requirement, not an implication of
      the finite betting test alone.
\item Bayesian confirmation and model comparison evaluate evidence
      under a stated prior/likelihood model
      (\cref{sec:p15-bayes-odds,sec:p15-bayesian-model-comparison}).
      Tight consequence here ranges over \emph{all} coherent extensions
      of the given premises and does not silently select a prior,
      update, or conditional-independence hypothesis.
\item The lower-prevision treatment has one-sided acceptable buying
      prices, full gamble domains, and separately specified natural or
      regular conditional extensions
      (\cref{def:p15-ip-coherence,def:p15-ip-conditional,def:p15-ip-updates}).
      Interval constraints here describe possible precise selections;
      equal-looking intervals are not automatically the same coherent
      lower-prevision assessment or update rule.
\item PCTL compares a path-event measure with a threshold in a Markov
      model, possibly optimizing over schedulers (\cref{def:p09-pctl}).
      It supplies neither this conditional-event semantics nor its
      Dutch-book consequence relation.
\end{itemize}

The chapter's final proposal combines qualitative nonmonotonic and
connexive principles with quantitative beliefs. Its suggested explanation
for the earlier neglect of nonclassical logics is research tradition,
not a demonstrated causal finding. Its cited technical results,
experimental reports, and future directions have been kept separate;
its reference list is not a new obligation to reconstruct every cited
book in full.
